
\begin{figure}[H]
  \centering
  \includegraphics[width=\linewidth]{parallelism_chips.pdf}
  \caption{What each JAX transformation does with the eight-chip slice.
  (a) Steady-state throughput on one chip with \texttt{jax.vmap} at
  batch sizes 1--8, and with \texttt{jax.pmap} mapping eight independent
  queries onto eight chips; the dashed line is the \texttt{vmap}
  batch-size-1 rate (2.051 proteins/second), not the 2.128
  proteins/second single-query baseline behind the 6.92$\times$ figure
  in the text, which comes from a different input family and timer.
  (b) HBM in use on each chip after the second call. Under the default
  path only one chip holds data; the auto-mesh summary reports 463\,MB
  per chip, consistent with replication, stored as a single scalar,
  with no per-device list, and drawn as a dashed line;
  \texttt{pmap} and the \texttt{pmap}+\texttt{pmean} ensemble
  distribute across all eight.
  Row labels give the results-catalogue row IDs.}
  \label{fig:par-chips}
\end{figure}

\subsection{Default execution leaves most of the slice idle}
\label{sec:par-default}

A Cloud TPU~v5e slice with eight chips exposes eight independent
devices to JAX, but nothing about calling \texttt{jax.jit} or invoking
AlphaFold2's own inference function causes computation to spread across
them. Under the default call path, only the first device (\texttt{TPU\_0})
allocates HBM for the model's parameters and activations; the
remaining seven report zero bytes used for the duration of the run.
The rest of the section covers two ways of asking JAX to use the other
seven, one that does not work and one that does, an automatic mechanism
that does not split the model either, and an ensemble built outside
the model.
Figure~\ref{fig:par-chips} summarizes all four cases.

\subsection{Batching within one device: \texttt{jax.vmap}}
\label{sec:par-vmap}

The first attempt batches queries with \texttt{jax.vmap}, which maps a
function over a leading array axis by tracing it once and vectorizing
the resulting operations. This is a within-device transformation: it
changes the shape of the computation that runs on a single core, not
the number of cores it runs on. Applied to AlphaFold2's forward pass, it
does not shard the batch axis across the slice. The sweep records only
\texttt{TPU\_0}'s HBM, which grows with batch size from 487\,MB to
821\,MB. The run write-up reports the other seven chips at 0\,MB, but no
per-device record was retained.

The throughput data confirms this is not a partial win. At batch sizes
1, 2, 4, and 8 the measured throughput is 2.051, 1.928, 1.487, and
1.508 proteins/second, i.e. 0.94$\times$, 0.73$\times$, and
0.74$\times$ the batch-size-1 rate.
No tested batch size exceeds the batch-size-1 throughput, and the
degradation is not monotonic: the batch-8 rate is marginally higher
than batch-4. We report the values as measured rather than fitting a
smooth trend to them. At every batch size tested, vectorization within
a device gives Cloud TPU~v5e no more throughput than not batching at
all.

\subsection{Explicit device mapping: \texttt{jax.pmap}}
\label{sec:par-pmap}

\texttt{jax.pmap} traces the function once and executes one copy per
device, communicating explicitly rather than relying on an
automatic partitioner. Assigning eight independent protein queries to
the eight chips, one query per chip, yields
14.718~proteins/second against a 2.128~proteins/second single-query,
single-chip baseline (0.47\,s per query): a 6.92$\times$ throughput
gain.
Distinct per-chip HBM allocations confirm the work is distributed
across the chips: 644\,MB on
\texttt{TPU\_0} and 445--469\,MB on the remaining seven chips.
On a like-for-like eight-versus-one-chip comparison at a fixed input
length of 100 residues, the corresponding speedup is 6.53$\times$,
that is 81.6\% parallel efficiency; the 118-residue multi-query run's
6.92$\times$ corresponds to 86.5\%.
The two figures answer different questions: 6.92$\times$ is the
throughput gain over a single query on a single chip, and 6.53$\times$
is the eight-over-one-chip speedup on a matched workload,
corresponding to 81.6\% parallel efficiency. Neither should be quoted
in place of the other.

The difference between Section~\ref{sec:par-vmap} and this section is one of
semantics, not tuning. \texttt{vmap} vectorizes an axis
\emph{within} a computation; on its own, with unsharded inputs as
here, it gives the partitioner no reason to place slices of that axis
on different devices. \texttt{pmap} does exactly that: it is, by construction, a
per-device dispatch.

\subsection{Automatic sharding}
\label{sec:par-gspmd}

Between manual \texttt{pmap} and no parallelism at all sits JAX's
automatic partitioner, which infers a sharding of a computation graph
from the sharding of its inputs and outputs without requiring the
programmer to rewrite the function. Two implementations of that idea
exist: GSPMD \citep{xu2021gspmd}, long the default, and Shardy
\citep{shardy2026}, which has since replaced it as JAX's default. The
committed Job configuration pins \texttt{jax[tpu]==0.10.2} and sets no
partitioner flag; that release defaults to Shardy. No retained
execution artifact records the effective setting, so we describe the
automatic-mesh experiment without attributing its outcome to a
confirmed implementation. Both implementations propagate shardings
outward from annotations, which is the property the observation below
turns on, though shared principles do not guarantee identical
decisions on a given graph.
Running AlphaFold2's forward pass under automatic mesh assignment
shows no sign of a sharded model: both retained run summaries record
463\,MB of HBM per chip, equal to the unsharded single-chip footprint,
with steady states of 0.472\,s and 0.473\,s. Each summary stores that
one per-chip figure, with no per-device list, and no compiled
layout or partitioner decision was retained.
That is consistent with replication and gives no evidence of any
reduction in per-chip footprint. It does not by itself establish that
every chip executed an identical full computation.

An earlier account of this result attributed the failure to AlphaFold2's
ensembling, implemented as a sequential \texttt{hk.while\_loop} with no
axis for the partitioner to distribute over. That explanation does not
hold for this experiment: the run in question used
\texttt{num\_ensemble\,=\,1}. AlphaFold2 computes the first ensemble
member outside the loop and enters \texttt{hk.while\_loop} at index~1
with the condition $i <$ \texttt{num\_ensemble}, so during inference
at a single member the loop body never executes (initialization
invokes it once, separately), and it cannot be the cause of the
single-chip-sized footprint the run records. AlphaFold2's released evaluation configuration in fact sets
\texttt{num\_ensemble\,=\,1} by default; the model does not run an
ensemble unless explicitly configured to.

The explanation we find most plausible does not depend on control
flow at all, and we state it as a hypothesis rather than a
demonstrated cause. An automatic partitioner of this kind works by
propagating sharding information outward from whatever annotations
exist on a computation's inputs and parameters. Inspecting AlphaFold2's
Haiku modules and our driver, we find no partition spec attached to
any weight or activation, and none supplied from outside through
input or output shardings. A partitioner given nothing to propagate
has no basis for a split, and replicating the whole computation is
its documented fallback \citep{xu2021gspmd}. DrJAX reports, by
analogy, that removing explicit sharding structure from a
partitioner-based system removes its ability to scale even for
embarrassingly parallel patterns \citep{drjax2024}. Our records do not
contain the evidence that would confirm this account for this graph:
no compiler trace, no retained array layouts, and no ablation that
adds annotations and observes a change. We present the missing
annotations as the likely explanation, consistent with the design
literature and with what we observed, not as a result we measured.

The sequential \texttt{hk.while\_loop} used for AlphaFold2's ensembling
does matter for the external ensemble of Section~\ref{sec:par-fix}:
because the ensemble average is computed inside
a loop with no vectorizable or shardable axis, distributing it
requires an approach that does not depend on the partitioner inferring structure
that the loop does not expose.

\subsection{An external ensemble with \texttt{pmap} + \texttt{lax.pmean}}
\label{sec:par-fix}

Rather than modify AlphaFold2's sequential loop, we sidestep it and build
an ensemble outside the model. The configuration keeps
\texttt{num\_ensemble\,=\,1}, so each chip runs exactly the
single-member forward pass AlphaFold2 always runs; what varies across
chips is the featurization, built eight times with a different
\texttt{random\_seed} each. An explicit \texttt{pmap} places one member
per chip and a \texttt{jax.lax.pmean} averages the resulting
\texttt{predicted\_lddt} logits across devices.
Listing~\ref{lst:fix} contrasts the two shapes.

This demonstration has two limits. AlphaFold2's own ensembling is not
replaced or exercised here: it stays
at a single member throughout, and the loop we describe above never
executes its body, so this experiment shows that \emph{an}
ensemble can be distributed across the slice, not that AlphaFold2's
internal one can. The quantity averaged also differs: AlphaFold2 combines
representations inside the model, whereas \texttt{pmean} here reduces
the per-member confidence logits after the forward pass, which is
where a plain JAX array is available to a collective. The eight members
share one trivial single-sequence MSA and differ only by seed, so they
are near-identical inputs; the result demonstrates the mechanism, not a
scientifically meaningful ensemble.

\noindent\begin{minipage}{\linewidth}
\begin{lstlisting}[caption={AlphaFold2's internal ensembling, which stays at a single member here, against the external ensemble we distribute instead. The per-member forward pass is the same in both.}, label={lst:fix}, language=Python]
# AlphaFold2's own ensembling: the first member is computed outside
# the loop, and hk.while_loop accumulates the rest sequentially on one
# device. We leave this untouched at num_ensemble = 1, so the loop body
# never executes.
representations = evoformer(slice_batch(0))

def body(x):
    i, current = x
    return i + 1, current + evoformer(slice_batch(i))

_, total = hk.while_loop(lambda x: x[0] < num_ensemble, body,
                         (1, representations))

# What we run instead: one differently seeded featurization per chip,
# each a standard single-member forward pass, averaged across devices.
per_member_features = [featurize(query, random_seed=i)
                       for i in range(n_chips)]

@functools.partial(jax.pmap, axis_name="ensemble")
def sharded_apply(feat):
    out = runner.apply(params, rng, feat)
    return jax.lax.pmean(out["predicted_lddt"]["logits"],
                         axis_name="ensemble")

averaged = sharded_apply(stack(per_member_features))
\end{lstlisting}
\end{minipage}

This places the work on the whole slice: all eight chips report nonzero
HBM allocation, and an \texttt{allclose} check on the first and last
returned replicas passes.
The change touches none of AlphaFold2's own modules; only the entry point
that drives ensemble members is replaced. \texttt{jax.lax.pmean}
averages the \texttt{predicted\_lddt} logits emitted by each member;
individual predicted structures are not themselves averaged.
Steady-state latency for the eight-way ensemble-sharded configuration
is 0.538\,s.
No sequential \texttt{num\_ensemble=8} baseline exists in the data, so
no speedup figure is reported or derived for this configuration; it
should not be described as faster or slower than a sequential
ensemble we did not measure.

This result is distinct from Section~\ref{sec:par-pmap}'s multi-query
experiment, which maps eight \emph{independent} protein queries onto
eight chips and measures throughput. Here the eight \emph{members of
one query's ensemble} are mapped onto eight chips, and the question is
whether the reduction that combines them can be executed without a
sequential loop. Both use \texttt{pmap}, but they parallelize different
axes of the problem.

\subsection{Summary}
\label{sec:par-summary}

The slice's eight chips are never disabled or defective. Every
experiment in this section runs on the same hardware. What changes
across the section is how much of that hardware the program
uses: one chip under the default path and under \texttt{vmap},
eight chips under \texttt{pmap} for independent queries, with a
6.92$\times$ throughput gain over a single query at 118 residues and a
6.53$\times$ eight-over-one-chip speedup at 100 residues, one chip's
worth of computation
apparently replicated eight times under automatic sharding, and
eight chips each running a full forward pass once an ensemble is built
outside the model and mapped across them with an explicit
\texttt{pmap}+\texttt{pmean} reduction. The last case distributes
independent members and reduces their confidence logits; it does not
partition one model instance, and does not repair the automatic
sharding of the case before it. In all four cases the hardware is the
same, and the outcome depends on the JAX transformation used to
express the work.
